\subsection{\label{sec.sf.schur}Schur polynomials}

\subsubsection{Alternants}

Here is one way to generate symmetric polynomials:

\begin{example}
Let $N=3$, and let us again abbreviate the indeterminates $x_{1},x_{2},x_{3}$
as $x,y,z$. For simplicity, we assume that $K$ is a field. As we know (from
the Vandermonde determinant -- specifically, Theorem \ref{thm.det.vander}
\textbf{(a)}), we have%
\[
\det\left(
\begin{array}
[c]{ccc}%
x^{2} & x & 1\\
y^{2} & y & 1\\
z^{2} & z & 1
\end{array}
\right)  =\prod_{i<j}\left(  x_{i}-x_{j}\right)  =\left(  x-y\right)  \left(
x-z\right)  \left(  y-z\right)  .
\]


What about similar determinants, such as $\det\left(
\begin{array}
[c]{ccc}%
x^{5} & x^{3} & 1\\
y^{5} & y^{3} & 1\\
z^{5} & z^{3} & 1
\end{array}
\right)  $ ? Just as in the proof of Lemma \ref{lem.det.vander.a.pol} (in
which we computed the original Vandermonde determinant), we can argue that
this is a polynomial in $x,y,z$ that is divisible by each of $x-y$ and $x-z$
and $y-z$ (since it becomes $0$ if we set one of $x,y,z$ equal to another).
Hence,%
\[
\det\left(
\begin{array}
[c]{ccc}%
x^{5} & x^{3} & 1\\
y^{5} & y^{3} & 1\\
z^{5} & z^{3} & 1
\end{array}
\right)  =\left(  x-y\right)  \left(  x-z\right)  \left(  y-z\right)  \cdot q
\]
for some $q\in K\left[  x,y,z\right]  $. However, this time, degree
considerations yield $\deg q=8-3=5$, so $q$ is no longer just a constant. What
is $q$ ? Using computer algebra, we see that%
\begin{align*}
q  &  =\dfrac{\det\left(
\begin{array}
[c]{ccc}%
x^{5} & x^{3} & 1\\
y^{5} & y^{3} & 1\\
z^{5} & z^{3} & 1
\end{array}
\right)  }{\left(  x-y\right)  \left(  x-z\right)  \left(  y-z\right)
}=\dfrac{-x^{5}y^{3}+x^{5}z^{3}+x^{3}y^{5}-x^{3}z^{5}-y^{5}z^{3}+y^{3}z^{5}%
}{-x^{2}y+x^{2}z+xy^{2}-xz^{2}-y^{2}z+yz^{2}}\\
&  =x^{2}y^{3}+x^{3}y^{2}+x^{2}z^{3}+x^{3}z^{2}+y^{2}z^{3}+y^{3}z^{2}%
+xyz^{3}\\
&  \ \ \ \ \ \ \ \ \ \ +xy^{3}z+x^{3}yz+2xy^{2}z^{2}+2x^{2}yz^{2}+2x^{2}%
y^{2}z\\
&  =m_{\left(  3,2,0\right)  }+m_{\left(  3,1,1\right)  }+2m_{\left(
2,2,1\right)  }\in\mathcal{S}.
\end{align*}
Note that $q\in\mathcal{S}$ can be easily seen without computing $q$. Indeed,
if we swap two of our variables $x,y,z$, then both $\det\left(
\begin{array}
[c]{ccc}%
x^{5} & x^{3} & 1\\
y^{5} & y^{3} & 1\\
z^{5} & z^{3} & 1
\end{array}
\right)  $ and $\left(  x-y\right)  \left(  x-z\right)  \left(  y-z\right)  $
get multiplied by $-1$, so their ratio $q$ stays unchanged. This shows that
$\sigma\cdot q=q$ whenever $\sigma\in S_{3}$ is a transposition. Since the
transpositions generate the group $S_{3}$ (indeed, Corollary
\ref{cor.perm.generated} yields that the simple transpositions $s_{1},s_{2}$
generate $S_{3}$), this entails that $\sigma\cdot q=q$ for any $\sigma\in
S_{3}$ (not just for transpositions). This means that $q$ is symmetric.

There is nothing special about the exponents $5$ and $3$ and $0$ in the above
determinant. More generally, for any $a,b,c\in\mathbb{N}$, we can define the
so-called \emph{alternant}%
\[
\det\left(
\begin{array}
[c]{ccc}%
x^{a} & x^{b} & x^{c}\\
y^{a} & y^{b} & y^{c}\\
z^{a} & z^{b} & z^{c}%
\end{array}
\right)  \in\mathcal{P}.
\]
When studying this alternant, we can WLOG assume that $a,b,c$ are distinct
(since otherwise, the alternant is just $0$) and furthermore assume that
$a>b>c$ (since the general case is reduced to this one by swapping the columns
around). The alternant is then a polynomial divisible by $x-y$ and $x-z$ and
$y-z$ (since it becomes $0$ if we set one of $x,y,z$ equal to another), and
thus divisible by $\left(  x-y\right)  \left(  x-z\right)  \left(  y-z\right)
=\det\left(
\begin{array}
[c]{ccc}%
x^{2} & x & 1\\
y^{2} & y & 1\\
z^{2} & z & 1
\end{array}
\right)  $ (the simplest nonzero alternant). Moreover, the ratio%
\[
\dfrac{\det\left(
\begin{array}
[c]{ccc}%
x^{a} & x^{b} & x^{c}\\
y^{a} & y^{b} & y^{c}\\
z^{a} & z^{b} & z^{c}%
\end{array}
\right)  }{\left(  x-y\right)  \left(  x-z\right)  \left(  y-z\right)
}=\dfrac{\det\left(
\begin{array}
[c]{ccc}%
x^{a} & x^{b} & x^{c}\\
y^{a} & y^{b} & y^{c}\\
z^{a} & z^{b} & z^{c}%
\end{array}
\right)  }{\det\left(
\begin{array}
[c]{ccc}%
x^{2} & x & 1\\
y^{2} & y & 1\\
z^{2} & z & 1
\end{array}
\right)  }%
\]
is a symmetric polynomial in $x,y,z$ (by the same argument that we used
before). Some experimentation suggests that all coefficients of this
polynomial are positive integers (to be rigorous, nonnegative integers). There
is probably no way of showing this without explicitly finding this polynomial
-- and the best way to do so is by defining this polynomial combinatorially.
\end{example}

Let us prepare for doing this.

\begin{definition}
\label{def.sf.alternants}\textbf{(a)} We let $\rho$ be the $N$-tuple $\left(
N-1,N-2,\ldots,N-N\right)  \in\mathbb{N}^{N}$. \medskip

\textbf{(b)} For any $N$-tuple $\alpha=\left(  \alpha_{1},\alpha_{2}%
,\ldots,\alpha_{N}\right)  \in\mathbb{N}^{N}$, we define%
\[
a_{\alpha}:=\det\left(  \underbrace{\left(  x_{i}^{\alpha_{j}}\right)  _{1\leq
i\leq N,\ 1\leq j\leq N}}_{\in\mathcal{P}^{N\times N}}\right)  \in
\mathcal{P}.
\]
This is called the $\alpha$\emph{-alternant} (of $x_{1},x_{2},\ldots,x_{N}$).
\end{definition}

For example, for $N=3$, we have%
\[
a_{\left(  5,3,0\right)  }=\det\left(
\begin{array}
[c]{ccc}%
x^{5} & x^{3} & 1\\
y^{5} & y^{3} & 1\\
z^{5} & z^{3} & 1
\end{array}
\right)  \ \ \ \ \ \ \ \ \ \ \left(  \text{where }\left(  x,y,z\right)
=\left(  x_{1},x_{2},x_{3}\right)  \right)  .
\]


Note that the definition of $a_{\rho}$ yields%
\begin{align}
a_{\rho}  &  =\det\left(  \left(  x_{i}^{\rho_{j}}\right)  _{1\leq i\leq
N,\ 1\leq j\leq N}\right)  =\det\left(  \left(  x_{i}^{N-j}\right)  _{1\leq
i\leq N,\ 1\leq j\leq N}\right) \nonumber\\
&  \ \ \ \ \ \ \ \ \ \ \ \ \ \ \ \ \ \ \ \ \left(  \text{since }\rho
_{j}=N-j\text{ for each }j\in\left[  N\right]  \right) \nonumber\\
&  =\prod_{1\leq i<j\leq N}\left(  x_{i}-x_{j}\right)
\label{eq.def.sf.alternants.arho=vdm}%
\end{align}
(by Theorem \ref{thm.det.vander} \textbf{(a)}, applied to $N$, $\mathcal{P}$
and $x_{i}$ instead of $n$, $K$ and $a_{i}$).

Thus, we suspect:

\begin{conjecture}
\label{conj.sf.schur-alt}For every $\alpha=\left(  \alpha_{1},\alpha
_{2},\ldots,\alpha_{N}\right)  \in\mathbb{N}^{N}$, the alternant $a_{\alpha}$
is a multiple of $a_{\rho}$ in the polynomial ring $\mathcal{P}$. Furthermore,
if $\alpha_{1}>\alpha_{2}>\cdots>\alpha_{N}$, then the polynomial $a_{\alpha
}/a_{\rho}$ has positive (more precisely, nonnegative) integer coefficients.
\end{conjecture}

We will prove this by explicitly constructing $a_{\alpha}/a_{\rho}$ combinatorially.

\subsubsection{Young diagrams and Schur polynomials}

Let us first define the \emph{Young diagram} of an $N$-partition. This is
analogous to the definition of the Young diagram of a partition (which we did
back in the proof of Proposition \ref{prop.pars.pkn=dual}):

\begin{definition}
\label{def.sf.ydiag}Let $\lambda$ be an $N$-partition.

The \emph{Young diagram} of $\lambda$ is defined to be a table of $N$
left-aligned rows, with the $i$-th row (counted from the top, as always)
having $\lambda_{i}$ boxes. Formally, the Young diagram of $\lambda$ is
defined as the set%
\[
\left\{  \left(  i,j\right)  \ \mid\ i\in\left[  N\right]  \text{ and }%
j\in\left[  \lambda_{i}\right]  \right\}  \subseteq\left\{  1,2,3,\ldots
\right\}  ^{2}.
\]
We visually represent each element $\left(  i,j\right)  $ of this Young
diagram as a box in row $i$ and column $j$; thus we obtain a table with $N$
left-aligned rows (some of which might be empty).

We denote the Young diagram of $\lambda$ by $Y\left(  \lambda\right)  $.
\end{definition}

For example, the $3$-partition $\left(  4,1,0\right)  $ has Young diagram%
\[
\ydiagram{4,1,0}\ \ .
\]
(The $3$-rd row is invisible since it has length $0$.) The four boxes in the
$1$-st (i.e., topmost) row of this diagram are $\left(  1,1\right)  $,
$\left(  1,2\right)  $, $\left(  1,3\right)  $ and $\left(  1,4\right)  $
(from left to right), while the single box in its $2$-nd row is $\left(
2,1\right)  $.

Now, we are going to fill our Young diagrams -- i.e., to put numbers in the boxes:

\begin{definition}
\label{def.sf.ytab}Let $\lambda$ be an $N$-partition.

A \emph{Young tableau} of shape $\lambda$ means a way of filling the boxes of
$Y\left(  \lambda\right)  $ with elements of $\left[  N\right]  $ (one element
per box). Formally speaking, it is defined as a map $T:Y\left(  \lambda
\right)  \rightarrow\left[  N\right]  $. We visually represent such a map by
filling in the number $T\left(  i,j\right)  $ into each box $\left(
i,j\right)  $.

We often abbreviate \textquotedblleft Young tableau\textquotedblright\ as
\textquotedblleft\emph{tableau}\textquotedblright. The plural of
\textquotedblleft tableau\textquotedblright\ is \textquotedblleft
tableaux\textquotedblright.
\end{definition}

For instance, here is a Young tableau of shape $\left(  4,3,3,0,0,0,0,\ldots
,0\right)  $ (defined for any $N\geq7$):%
\[
\ytableaushort{1724,336,261}\ \ .
\]
Formally speaking, this is a map $T:Y\left(  4,3,3,0,0,0,0,\ldots,0\right)
\rightarrow\left[  N\right]  $ that sends the pairs $\left(  1,1\right)
,\left(  1,2\right)  ,\left(  1,3\right)  ,\left(  1,4\right)  ,\left(
2,1\right)  ,\ldots$ to $1,7,2,4,3,\ldots$, respectively.

We will use some visually inspired language when talking about Young diagrams
and tableaux:

\begin{itemize}
\item For instance, the \emph{entry} of a tableau $T$ in box $\left(
i,j\right)  $ will mean the value $T\left(  i,j\right)  $.

\item Also, the $u$\emph{-th row} of a tableau $T$ (for a given $u\geq1$) will
mean the sequence of all entries of $T$ in the boxes $\left(  i,j\right)  $
with $i=u$.

\item Likewise, the $v$\emph{-th column} of a tableau $T$ (for a given
$v\geq1$) will mean the sequence of all entries of $T$ in the boxes $\left(
i,j\right)  $ with $j=v$.

\item If $T$ is a Young tableau of shape $\lambda$, then the boxes of
$Y\left(  \lambda\right)  $ will also be called the boxes of $T$.

\item Two boxes of a Young diagram (or of a tableau) are said to be
\emph{adjacent} if they have an edge in common when drawn on the picture
(i.e., when one of them has the form $\left(  i,j\right)  $, while the other
has the form $\left(  i,j+1\right)  $ or $\left(  i+1,j\right)  $).

\item The words \textquotedblleft north\textquotedblright, \textquotedblleft
west\textquotedblright, \textquotedblleft south\textquotedblright\ and
\textquotedblleft east\textquotedblright\ are to be understood according to
the picture of a Young diagram: e.g., the box $\left(  2,4\right)  $ lies one
step north and three steps west of the box $\left(  3,7\right)  $.
\end{itemize}

Some tableaux are better than others:

\begin{definition}
\label{def.sf.ssyt}Let $\lambda$ be an $N$-partition.

A Young tableau $T$ of shape $\lambda$ is said to be \emph{semistandard} if
its entries

\begin{itemize}
\item increase weakly along each row (from left to right);

\item increase strictly down each column (from top to bottom).
\end{itemize}

Formally speaking, this means that a Young tableau $T:Y\left(  \lambda\right)
\rightarrow\left[  N\right]  $ is semistandard if and only if

\begin{itemize}
\item we have $T\left(  i,j\right)  \leq T\left(  i,j+1\right)  $ for any
$\left(  i,j\right)  \in Y\left(  \lambda\right)  $ satisfying $\left(
i,j+1\right)  \in Y\left(  \lambda\right)  $;

\item we have $T\left(  i,j\right)  <T\left(  i+1,j\right)  $ for any $\left(
i,j\right)  \in Y\left(  \lambda\right)  $ satisfying $\left(  i+1,j\right)
\in Y\left(  \lambda\right)  $.
\end{itemize}

We let $\operatorname*{SSYT}\left(  \lambda\right)  $ denote the set of all
semistandard Young tableaux of shape $\lambda$. (This depends on $N$ as well,
but $N$ is fixed, so we omit it from our notation.) We will usually say
\textquotedblleft\emph{semistandard tableau}\textquotedblright\ instead of
\textquotedblleft semistandard Young tableau\textquotedblright.
\end{definition}

\begin{example}
Consider the following $6$ Young tableaux of shape $\left(  4,2,1,0,0,0,\ldots
,0\right)  $:%
\begin{align}
&  \ytableaushort{1334,235,4}\qquad\ytableaushort{2134,345,6}\qquad
\ytableaushort{1123,245,6}\label{eq.def.sf.ssyt.1.1}\\
& \nonumber\\
&  \ytableaushort{1234,567,8}\qquad\ytableaushort{1111,222,3}\qquad
\ytableaushort{1234,123,1}\nonumber
\end{align}
Which of these $6$ tableaux are semistandard? The first one is not
semistandard, since the entries in its second column do not strictly increase
down the column. The second one is not semistandard, since the entries in its
first row do not weakly increase along the row. The third one is semistandard.
The fourth one is semistandard, too. The fifth one is semistandard, too
(actually it has a special property: each of its entries is the smallest
possible value that an entry of a semistandard tableau could have in its box).
The sixth one is not semistandard, again because of the columns.
\end{example}

\begin{definition}
\label{def.sf.ytab.xT}Let $\lambda$ be an $N$-partition. If $T$ is any Young
tableau of shape $\lambda$, then we define the corresponding monomial%
\[
x_{T}:=\prod_{c\text{ is a box of }Y\left(  \lambda\right)  }x_{T\left(
c\right)  }=\prod_{\left(  i,j\right)  \in Y\left(  \lambda\right)
}x_{T\left(  i,j\right)  }=\prod_{k=1}^{N}x_{k}^{\left(  \text{\# of times
}k\text{ appears in }T\right)  }.
\]

\end{definition}

For example, the three Young tableaux in (\ref{eq.def.sf.ssyt.1.1}) have
corresponding monomials%
\begin{align*}
&  x_{1}x_{3}x_{3}x_{4}x_{2}x_{3}x_{5}x_{4}=x_{1}x_{2}x_{3}^{3}x_{4}^{2}%
x_{5},\\
&  x_{2}x_{1}x_{3}x_{4}x_{3}x_{4}x_{5}x_{6},\\
&  x_{1}x_{1}x_{2}x_{3}x_{2}x_{4}x_{5}x_{6}.
\end{align*}


\begin{definition}
\label{def.sf.schur}Let $\lambda$ be an $N$-partition. We define the
\emph{Schur polynomial} $s_{\lambda}\in\mathcal{P}$ by%
\[
s_{\lambda}:=\sum_{T\in\operatorname*{SSYT}\left(  \lambda\right)  }x_{T}.
\]

\end{definition}

\begin{example}
\label{exa.sf.schur-h-e}\textbf{(a)} Let $n\in\mathbb{N}$. Consider the
$N$-partition $\left(  n,0,0,\ldots,0\right)  $. The semistandard tableaux $T$
of shape $\left(  n,0,0,\ldots,0\right)  $ are simply the fillings of a single
row with $n$ elements of $\left[  N\right]  $ that weakly increase from left
to right:%
\[
T=\begin{ytableau}
i_1 & i_2 & \cdots & i_n
\end{ytableau}\ \ \ \ \ \ \ \ \ \ \text{with }i_{1}\leq i_{2}\leq\cdots\leq
i_{n}.
\]
Thus,
\[
s_{\left(  n,0,0,\ldots,0\right)  }=\sum_{i_{1}\leq i_{2}\leq\cdots\leq i_{n}%
}x_{i_{1}}x_{i_{2}}\cdots x_{i_{n}}=h_{n}.
\]


\textbf{(b)} Let $n\in\left\{  0,1,\ldots,N\right\}  $. Consider the
$N$-partition $\left(  1,1,\ldots,1,0,0,\ldots,0\right)  $ (with $n$ ones and
$N-n$ zeroes). The semistandard tableaux $T$ of shape $\left(  1,1,\ldots
,1,0,0,\ldots,0\right)  $ are simply the fillings of a single column with $n$
elements of $\left[  N\right]  $ that strictly increase from top to bottom:%
\[
T=\begin{ytableau}
i_1 \\ i_2 \\ \vdots \\ i_n
\end{ytableau}\ \ \ \ \ \ \ \ \ \ \text{with }i_{1}<i_{2}<\cdots<i_{n}.
\]
Hence,%
\[
s_{\left(  1,1,\ldots,1,0,0,\ldots,0\right)  \text{ (with }n\text{ ones and
}N-n\text{ zeroes)}}=\sum_{i_{1}<i_{2}<\cdots<i_{n}}x_{i_{1}}x_{i_{2}}\cdots
x_{i_{n}}=e_{n}.
\]


\textbf{(c)} Assume that $N\geq2$. Consider the $N$-partition $\left(
2,1,0,0,0,\ldots,0\right)  $ (with all entries from the third on being $0$).
The semistandard tableaux $T$ of shape $\left(  2,1,0,0,0,\ldots,0\right)  $
all have the form%
\[
T=\ytableaushort{ij,k}\ \ \ \ \ \ \ \ \ \ \text{with }i\leq j\text{ and }i<k.
\]
Hence,%
\begin{align*}
s_{\left(  2,1,0,0,0,\ldots,0\right)  }  &  =\sum_{\substack{i\leq
j;\\i<k}}x_{i}x_{j}x_{k}=\underbrace{\sum_{\substack{i<k;\\j=i}}x_{i}%
x_{j}x_{k}}_{=\sum_{i<k}x_{i}x_{i}x_{k}}+\underbrace{\sum_{i<j<k}x_{i}%
x_{j}x_{k}}_{=e_{3}}+\underbrace{\sum_{\substack{i<k;\\j=k}}x_{i}x_{j}x_{k}%
}_{=\sum_{i<k}x_{i}x_{k}x_{k}}+\underbrace{\sum_{i<k<j}x_{i}x_{j}x_{k}%
}_{=e_{3}}\\
&  \ \ \ \ \ \ \ \ \ \ \ \ \ \ \ \ \ \ \ \ \left(
\begin{array}
[c]{c}%
\text{since each triple }\left(  i,j,k\right)  \text{ of elements of }\left[
N\right] \\
\text{that satisfies }i\leq j\text{ and }i<k\text{ must satisfy}\\
\text{\textbf{exactly one} of the four}\\
\text{conditions }\left(  i<k\text{ and }j=i\right)  \text{ and }i<j<k\\
\text{and }\left(  i<k\text{ and }j=k\right)  \text{ and }i<k<j\text{,}\\
\text{and conversely, each triple satisfying one}\\
\text{of the latter four conditions must}\\
\text{satisfy }i\leq j\text{ and }i<k
\end{array}
\right) \\
&  =\sum_{i<k}x_{i}x_{i}x_{k}+e_{3}+\sum_{i<k}x_{i}x_{k}x_{k}+e_{3}%
=2e_{3}+\sum_{i<k}\underbrace{x_{i}x_{i}}_{=x_{i}^{2}}x_{k}+\sum_{i<k}%
x_{i}\underbrace{x_{k}x_{k}}_{=x_{k}^{2}}\\
&  =2e_{3}+\underbrace{\sum_{i<k}x_{i}^{2}x_{k}+\sum_{i<k}x_{i}x_{k}^{2}%
}_{=m_{\left(  2,1,0,0,\ldots,0\right)  }}=2e_{3}+\underbrace{m_{\left(
2,1,0,0,\ldots,0\right)  }}_{\substack{=e_{2}e_{1}-3e_{3}\\\text{(check
this!)}}}\\
&  =2e_{3}+e_{2}e_{1}-3e_{3}=e_{2}e_{1}-e_{3}.
\end{align*}

\end{example}

\begin{theorem}
\label{thm.sf.schur-symm}Let $\lambda$ be an $N$-partition. Then:\medskip

\textbf{(a)} The polynomial $s_{\lambda}$ is symmetric. \medskip

\textbf{(b)} We have
\[
a_{\lambda+\rho}=a_{\rho}\cdot s_{\lambda}.
\]
Here, the addition on $\mathbb{N}^{N}$ is defined entrywise: that is, if
$\alpha=\left(  \alpha_{1},\alpha_{2},\ldots,\alpha_{N}\right)  $ and
$\beta=\left(  \beta_{1},\beta_{2},\ldots,\beta_{N}\right)  $ are two
$N$-tuples in $\mathbb{N}^{N}$, then we set%
\[
\alpha+\beta:=\left(  \alpha_{1}+\beta_{1},\alpha_{2}+\beta_{2},\ldots
,\alpha_{N}+\beta_{N}\right)  .
\]

\end{theorem}

This theorem (once it will be proved) will yield Conjecture
\ref{conj.sf.schur-alt} in the case when $\alpha_{1}>\alpha_{2}>\cdots
>\alpha_{N}$. Indeed, if $\alpha=\left(  \alpha_{1},\alpha_{2},\ldots
,\alpha_{N}\right)  \in\mathbb{N}^{N}$ is an $N$-tuple satisfying $\alpha
_{1}>\alpha_{2}>\cdots>\alpha_{N}$, then we can write $\alpha$ as
$\alpha=\lambda+\rho$ for some $N$-partition $\lambda$ (namely, for
$\lambda=\alpha-\rho=\left(  \alpha_{1}-\left(  N-1\right)  ,\alpha
_{2}-\left(  N-2\right)  ,\ldots,\alpha_{N}-\left(  N-N\right)  \right)  $),
and then Theorem \ref{thm.sf.schur-symm} \textbf{(b)} will yield $a_{\alpha
}=a_{\rho}\cdot s_{\lambda}$, so that $a_{\alpha}/a_{\rho}=s_{\lambda}$, which
is a symmetric polynomial (by Theorem \ref{thm.sf.schur-symm} \textbf{(a)})
and furthermore has positive coefficients (by its combinatorial definition).
Once Conjecture \ref{conj.sf.schur-alt} is proved in the case when $\alpha
_{1}>\alpha_{2}>\cdots>\alpha_{N}$, the validity of its first claim in the
general case easily follows (because the alternant $a_{\alpha}$ is $0$ when
two of $\alpha_{1},\alpha_{2},\ldots,\alpha_{N}$ are equal\footnote{This is
easy to prove (but see Lemma \ref{lem.sf.alternant-0} \textbf{(a)} below for a
proof).}, and otherwise can be reduced to the case $\alpha_{1}>\alpha
_{2}>\cdots>\alpha_{N}$ by swapping the columns around\footnote{See Lemma
\ref{lem.sf.alternant-0} \textbf{(b)} below for the details of this.}).

We will prove Theorem \ref{thm.sf.schur-symm} \textbf{(a)} soon and Theorem
\ref{thm.sf.schur-symm} \textbf{(b)} later.

\subsubsection{Skew Young diagrams and skew Schur polynomials}

Before we get to the proof, let us generalize the situation a bit:

\begin{definition}
\label{def.sf.par-subset}Let $\lambda$ and $\mu$ be two $N$-partitions.

We say that $\mu\subseteq\lambda$ if and only if $Y\left(  \mu\right)
\subseteq Y\left(  \lambda\right)  $. Equivalently, $\mu\subseteq\lambda$ if
and only if%
\[
\text{each }i\in\left[  N\right]  \text{ satisfies }\mu_{i}\leq\lambda_{i}%
\]
(where we write $\mu$ and $\lambda$ as $\mu=\left(  \mu_{1},\mu_{2},\ldots
,\mu_{N}\right)  $ and $\lambda=\left(  \lambda_{1},\lambda_{2},\ldots
,\lambda_{N}\right)  $). Thus we have defined a partial order $\subseteq$ on
the set of all $N$-partitions.
\end{definition}

For example, we have $\left(  3,2,0\right)  \subseteq\left(  4,2,1\right)  $
(since $3\leq4$ and $2\leq2$ and $0\leq1$), but we don't have $\left(
2,2,0\right)  \subseteq\left(  3,1,0\right)  $ (since $2>1$). We can see this
on the Young diagrams:%
\begin{align*}
&  \text{we have }\ydiagram{3,2,0} \subseteq\ydiagram{4,2,1}\text{\ \ ,}\\
& \\
&  \text{but we don't have }\ydiagram{2,2,0} \subseteq
\ydiagram{3,1,0}\text{\ \ .}%
\end{align*}


\begin{definition}
\label{def.sf.skew-diag}Let $\lambda$ and $\mu$ be two $N$-partitions such
that $\mu\subseteq\lambda$. Then, we define the \emph{skew Young diagram}
$Y\left(  \lambda/\mu\right)  $ to be the set difference%
\begin{align*}
Y\left(  \lambda\right)  \setminus Y\left(  \mu\right)   &  =\left\{  \left(
i,j\right)  \ \mid\ i\in\left[  N\right]  \text{ and }j\in\left[  \lambda
_{i}\right]  \setminus\left[  \mu_{i}\right]  \right\} \\
&  =\left\{  \left(  i,j\right)  \ \mid\ i\in\left[  N\right]  \text{ and
}j\in\mathbb{Z}\text{ and }\mu_{i}<j\leq\lambda_{i}\right\}  .
\end{align*}

\end{definition}

For example,%
\[
Y\left(  \left(  4,3,1\right)  /\left(  2,1,0\right)  \right)
=\ydiagram{2+2,1+2,0+1}\ \ .
\]
We note that any row or column of a skew Young diagram $Y\left(  \lambda
/\mu\right)  $ is contiguous, i.e., has no holes between boxes. Better yet, if
$\left(  a,b\right)  $ and $\left(  e,f\right)  $ are two boxes of $Y\left(
\lambda/\mu\right)  $, then any box $\left(  c,d\right)  $ that lies
\textquotedblleft between\textquotedblright\ them (i.e., that satisfies $a\leq
c\leq e$ and $b\leq d\leq f$) must also belong to $Y\left(  \lambda
/\mu\right)  $. Let us state this as a lemma:

\begin{lemma}
[Convexity of skew Young diagrams]\label{lem.sf.skew-diag.convexity}Let
$\lambda$ and $\mu$ be two $N$-partitions such that $\mu\subseteq\lambda$. Let
$\left(  a,b\right)  $ and $\left(  e,f\right)  $ be two elements of $Y\left(
\lambda/\mu\right)  $. Let $\left(  c,d\right)  \in\mathbb{Z}^{2}$ satisfy
$a\leq c\leq e$ and $b\leq d\leq f$. Then, $\left(  c,d\right)  \in Y\left(
\lambda/\mu\right)  $.
\end{lemma}

\begin{proof}
[Hints to the proof of Lemma \ref{lem.sf.skew-diag.convexity}.]This follows
easily from the definition of $Y\left(  \lambda/\mu\right)  $ using the fact
that $\lambda$ and $\mu$ are weakly decreasing $N$-tuples. A detailed proof
can be found in Section \ref{sec.details.sf.schur}.
\end{proof}

Lemma \ref{lem.sf.skew-diag.convexity} is known as the \emph{convexity} of
$Y\left(  \lambda/\mu\right)  $ (albeit in a very specific sense of the word
\textquotedblleft convexity\textquotedblright).

Next, we can define Young tableaux of shape $\lambda/\mu$ whenever $\lambda$
and $\mu$ are two $N$-partitions satisfying $\mu\subseteq\lambda$. The
definition is analogous to Definition \ref{def.sf.ytab}, except that we are
only filling the boxes of $Y\left(  \lambda/\mu\right)  $ (rather than all
boxes of $Y\left(  \lambda\right)  $) this time:

\begin{definition}
\label{def.sf.skew-tab}Let $\lambda$ and $\mu$ be two $N$-partitions such that
$\mu\subseteq\lambda$. A \emph{Young tableau} of shape $\lambda/\mu$ means a
way of filling the boxes of $Y\left(  \lambda/\mu\right)  $ with elements of
$\left[  N\right]  $ (one element per box). Formally speaking, it is defined
as a map $T:Y\left(  \lambda/\mu\right)  \rightarrow\left[  N\right]  $. We
visually represent such a map by filling in the number $T\left(  i,j\right)  $
into each box $\left(  i,j\right)  $.

Young tableaux of shape $\lambda/\mu$ are often called \emph{skew Young
tableaux}.

If we don't have $\mu\subseteq\lambda$, then we agree that there are no Young
tableaux of shape $\lambda/\mu$.
\end{definition}

The notion of a semistandard tableau of shape $\lambda/\mu$ is, again, defined
in the same way as for shape $\lambda$:

\begin{definition}
\label{def.sf.skew-ssyt}Let $\lambda$ and $\mu$ be two $N$-partitions.

A Young tableau $T$ of shape $\lambda/\mu$ is said to be \emph{semistandard}
if its entries

\begin{itemize}
\item increase weakly along each row (from left to right);

\item increase strictly down each column (from top to bottom).
\end{itemize}

Formally speaking, this means that a Young tableau $T:Y\left(  \lambda
/\mu\right)  \rightarrow\left[  N\right]  $ is semistandard if and only if

\begin{itemize}
\item we have $T\left(  i,j\right)  \leq T\left(  i,j+1\right)  $ for any
$\left(  i,j\right)  \in Y\left(  \lambda/\mu\right)  $ satisfying $\left(
i,j+1\right)  \in Y\left(  \lambda/\mu\right)  $;

\item we have $T\left(  i,j\right)  <T\left(  i+1,j\right)  $ for any $\left(
i,j\right)  \in Y\left(  \lambda/\mu\right)  $ satisfying $\left(
i+1,j\right)  \in Y\left(  \lambda/\mu\right)  $.
\end{itemize}

We let $\operatorname*{SSYT}\left(  \lambda/\mu\right)  $ denote the set of
all semistandard Young tableaux of shape $\lambda/\mu$. We will usually say
\textquotedblleft\emph{semistandard tableau}\textquotedblright\ instead of
\textquotedblleft semistandard Young tableau\textquotedblright.
\end{definition}

For example, here is a semistandard Young tableau of shape $\left(
4,3,1\right)  /\left(  2,1,0\right)  $:%
\[
\ytableaushort{\none\none 13,\none 22,1}\ \ .
\]
Meanwhile, there are no Young tableaux of shape $\left(  3,2,1\right)
/\left(  2,2,2\right)  $ (since we don't have $\left(  2,2,2\right)
\subseteq\left(  3,2,1\right)  $), and thus the set $\operatorname*{SSYT}%
\left(  \left(  3,2,1\right)  /\left(  2,2,2\right)  \right)  $ is empty.

The phrases \textquotedblleft increase weakly along each row\textquotedblright%
\ and \textquotedblleft increase strictly down each column\textquotedblright%
\ in Definition \ref{def.sf.skew-diag} have been formalized in terms of
adjacent entries: e.g., we have declared \textquotedblleft increase weakly
along each row\textquotedblright\ to mean \textquotedblleft$T\left(
i,j\right)  \leq T\left(  i,j+1\right)  $\textquotedblright\ rather than
\textquotedblleft$T\left(  i,j_{1}\right)  \leq T\left(  i,j_{2}\right)  $
whenever $j_{1}\leq j_{2}$\textquotedblright. However, since any row or column
of $Y\left(  \lambda/\mu\right)  $ is contiguous, the latter stronger meaning
actually follows from the former. To wit:

\begin{lemma}
\label{lem.sf.skew-ssyt.increase}Let $\lambda$ and $\mu$ be two $N$%
-partitions. Let $T$ be a semistandard Young tableau of shape $\lambda/\mu$.
Then: \medskip

\textbf{(a)} If $\left(  i,j_{1}\right)  $ and $\left(  i,j_{2}\right)  $ are
two elements of $Y\left(  \lambda/\mu\right)  $ satisfying $j_{1}\leq j_{2}$,
then $T\left(  i,j_{1}\right)  \leq T\left(  i,j_{2}\right)  $. \medskip

\textbf{(b)} If $\left(  i_{1},j\right)  $ and $\left(  i_{2},j\right)  $ are
two elements of $Y\left(  \lambda/\mu\right)  $ satisfying $i_{1}\leq i_{2}$,
then $T\left(  i_{1},j\right)  \leq T\left(  i_{2},j\right)  $. \medskip

\textbf{(c)} If $\left(  i_{1},j\right)  $ and $\left(  i_{2},j\right)  $ are
two elements of $Y\left(  \lambda/\mu\right)  $ satisfying $i_{1}<i_{2}$, then
$T\left(  i_{1},j\right)  <T\left(  i_{2},j\right)  $. \medskip

\textbf{(d)} If $\left(  i_{1},j_{1}\right)  $ and $\left(  i_{2}%
,j_{2}\right)  $ are two elements of $Y\left(  \lambda/\mu\right)  $
satisfying $i_{1}\leq i_{2}$ and $j_{1}\leq j_{2}$, then $T\left(  i_{1}%
,j_{1}\right)  \leq T\left(  i_{2},j_{2}\right)  $. \medskip

\textbf{(e)} If $\left(  i_{1},j_{1}\right)  $ and $\left(  i_{2}%
,j_{2}\right)  $ are two elements of $Y\left(  \lambda/\mu\right)  $
satisfying $i_{1}<i_{2}$ and $j_{1}\leq j_{2}$, then $T\left(  i_{1}%
,j_{1}\right)  <T\left(  i_{2},j_{2}\right)  $.
\end{lemma}

\begin{proof}
[Hints to the proof of Lemma \ref{lem.sf.skew-ssyt.increase}.]This is easy
using Lemma \ref{lem.sf.skew-diag.convexity}. A detailed proof can be found in
Section \ref{sec.details.sf.schur}.
\end{proof}

We extend Definition \ref{def.sf.ytab.xT} to skew tableaux:

\begin{definition}
\label{def.sf.ytab.skew-xT}Let $\lambda$ and $\mu$ be two $N$-partitions. If
$T$ is any Young tableau of shape $\lambda/\mu$, then we define the
corresponding monomial%
\[
x_{T}:=\prod_{c\text{ is a box of }Y\left(  \lambda/\mu\right)  }x_{T\left(
c\right)  }=\prod_{\left(  i,j\right)  \in Y\left(  \lambda/\mu\right)
}x_{T\left(  i,j\right)  }=\prod_{k=1}^{N}x_{k}^{\left(  \text{\# of times
}k\text{ appears in }T\right)  }.
\]

\end{definition}

For example,%
\[
\text{if }T=\ytableaushort{\none\none 13,\none 22,344}\ \ \text{, then }%
x_{T}=x_{1}x_{3}x_{2}x_{2}x_{3}x_{4}x_{4}=x_{1}x_{2}^{2}x_{3}^{2}x_{4}^{2}.
\]


Finally, we generalize Definition \ref{def.sf.schur} to $\lambda/\mu$:

\begin{definition}
\label{def.sf.skew-schur}Let $\lambda$ and $\mu$ be two $N$-partitions. We
define the \emph{skew Schur polynomial} $s_{\lambda/\mu}\in\mathcal{P}$ by%
\[
s_{\lambda/\mu}:=\sum_{T\in\operatorname*{SSYT}\left(  \lambda/\mu\right)
}x_{T}.
\]

\end{definition}

\begin{example}
\textbf{(a)} We have%
\[
s_{\left(  3,3,2,0,0,0,\ldots,0\right)  /\left(  2,1,0,0,0,\ldots,0\right)
}=\sum_{i<j\geq k<\ell\geq m}x_{i}x_{j}x_{k}x_{\ell}x_{m},
\]
since the semistandard tableaux of shape $\left(  3,3,2,0,0,0,\ldots,0\right)
/\left(  2,1,0,0,0,\ldots,0\right)  $ have the form
$\ytableaushort{\none\none i,\none kj,m\ell}$ for $i,j,k,\ell,m\in\left[
N\right]  $ satisfying $i<j\geq k<\ell\geq m$. \medskip

\textbf{(b)} We have%
\[
s_{\left(  3,2,1,0,0,0,\ldots,0\right)  /\left(  2,1,0,0,0,\ldots,0\right)
}=\sum_{i,j,k}x_{i}x_{j}x_{k},
\]
since the semistandard tableaux of shape $\left(  3,2,1,0,0,0,\ldots,0\right)
/\left(  2,1,0,0,0,\ldots,0\right)  $ have the form
$\ytableaushort{\none\none i,\none j,k}$ for $i,j,k\in\left[  N\right]  $
satisfying no special requirements (as there are never two entries in the same
row or in the same column of the tableau). Thus,%
\[
s_{\left(  3,2,1,0,0,0,\ldots,0\right)  /\left(  2,1,0,0,0,\ldots,0\right)
}=\sum_{i,j,k}x_{i}x_{j}x_{k}=\left(  \sum_{i}x_{i}\right)  ^{3}=\left(
x_{1}+x_{2}+\cdots+x_{N}\right)  ^{3}.
\]

\end{example}

\begin{theorem}
\label{thm.sf.skew-schur-symm}Let $\lambda$ and $\mu$ be any two
$N$-partitions. Then, the polynomial $s_{\lambda/\mu}$ is symmetric.
\end{theorem}

\begin{remark}
\label{rmk.sf.skew-0}If we set $\mathbf{0}=\left(  0,0,\ldots,0\right)
\in\mathbb{N}^{N}$, then $s_{\lambda/\mathbf{0}}=s_{\lambda}$ (since the
semistandard tableaux of shape $\lambda/\mathbf{0}$ are precisely the
semistandard tableaux of shape $\lambda$). Hence, Theorem
\ref{thm.sf.skew-schur-symm} generalizes Theorem \ref{thm.sf.schur-symm}
\textbf{(a)}.
\end{remark}

We will now prove Theorem \ref{thm.sf.skew-schur-symm} bijectively, using a
beautiful set of combinatorial bijections known as the \emph{Bender--Knuth
involutions}.

